\subsection{Off-chain Capacity Signals}
\label{sec:offchain}

The commit-reveal scheme of \Cref{sec:ewma} provides MEV-resistant capacity
updates but incurs latency of at least two block intervals plus a commit timeout
(approximately 40 seconds on Base L2). For real-time routing decisions, this
is prohibitive. We introduce an off-chain signed attestation protocol that
reduces update latency to seconds while preserving the same on-chain
EWMA smoothing and security guarantees.

\paragraph{Protocol.}
Each sink $k$ periodically measures its available capacity and produces a
signed \emph{capacity attestation}:

\begin{definition}[Capacity Attestation]
A capacity attestation is an EIP-712 typed data structure
$a = (\tasktype, k, c, t_s, n)$ where $\tasktype$ is the task type,
$k$ is the sink address, $c$ is the declared capacity, $t_s$ is the
timestamp, and $n$ is a monotonically increasing nonce.
The sink signs $a$ as $\sigma_k = \mathrm{Sign}_{sk_k}(a)$.
\end{definition}

Signed attestations are disseminated via any off-chain channel (direct
API calls, P2P gossip, or posted to a relay). Any party may collect a batch
of attestations and submit them to the on-chain \texttt{OffchainAggregator}
contract, which performs the following for each attestation in the batch:

\begin{enumerate}
    \item Verify the ECDSA signature against the claimed sink address.
    \item Check timestamp freshness: $|t_{\text{now}} - t_s| < \Delta$
        (we set $\Delta = 600$ seconds).
    \item Check nonce monotonicity: $n > n_{\text{last}}(k)$ for the sink.
    \item If valid, call the \texttt{CapacityRegistry} to apply the same EWMA
        smoothing as the commit-reveal path (\Cref{sec:ewma}).
\end{enumerate}

Invalid attestations (bad signatures, stale timestamps, replayed nonces) are
silently rejected with logged events; valid ones update the smoothed capacity.

\paragraph{Latency comparison.}
Under the commit-reveal path, a capacity update requires a commit transaction
(1 block), a waiting period (20 blocks $\approx 40$s), and a reveal transaction
(1 block), totaling approximately 44 seconds. The off-chain path requires only
the time to collect and relay an attestation batch, plus one on-chain transaction
for the batch submission. In practice, this reduces end-to-end latency to
under 5 seconds on Base L2.

\paragraph{Gas comparison.}
We measure gas costs on Base Sepolia (see \Cref{sec:onchain-eval} for full benchmarks):

\begin{center}
\begin{tabular}{lrr}
\toprule
\textbf{Path} & \textbf{Gas} & \textbf{Latency} \\
\midrule
Commit + Reveal (on-chain) & 58{,}287 & $\sim$40s \\
Aggregated attestation (off-chain) & 9{,}595 & $<$5s \\
\midrule
\textbf{Reduction} & \textbf{83.5\%} & \textbf{$\sim$87.5\%} \\
\bottomrule
\end{tabular}
\end{center}

The 83.5\% gas reduction comes from eliminating the commit phase entirely:
the off-chain path skips hash storage, timeout checking, and the two-transaction
round trip. The EWMA update itself is identical in both paths.

\paragraph{Security equivalence.}
The off-chain attestation path provides the same steady-state security
guarantees as commit-reveal:

\begin{proposition}[Off-chain Security Equivalence]
\label{prop:offchain-security}
Under the off-chain attestation protocol with EWMA smoothing, the capacity
manipulation bounds from \Cref{sec:ewma} hold: a single attestation can
shift the smoothed capacity by at most $\alpha$ (30\%) of the reported change.
Stake requirements and slashing conditions are unchanged.
\end{proposition}

\begin{proof}
The EWMA update applied by the aggregator is identical to the reveal-phase
update: $\Csmooth(k, t+1) = \alpha \cdot c + (1-\alpha) \cdot \Csmooth(k, t)$.
Since the smoothing filter is the same, the maximum per-update impact is
$\alpha \cdot |c - \Csmooth(k, t)|$, bounded by the stake-derived capacity
cap $\text{cap}(S_k)$ from \Cref{sec:ewma}. The nonce monotonicity check
prevents replay, serving the same role as the commit hash.
\end{proof}

The trade-off is MEV resistance: commit-reveal hides capacity values during
the commit phase, while off-chain attestations are visible to anyone who
receives them. In practice, EWMA smoothing limits the exploitable value of
front-running a single capacity update, making this trade-off acceptable
for the significant latency and gas improvements.
